% Options for packages loaded elsewhere
\PassOptionsToPackage{unicode}{hyperref}
\PassOptionsToPackage{hyphens}{url}
%
\documentclass[
]{article}
\usepackage{amsmath,amssymb}
\usepackage{iftex}
\ifPDFTeX
  \usepackage[T1]{fontenc}
  \usepackage[utf8]{inputenc}
  \usepackage{textcomp} % provide euro and other symbols
\else % if luatex or xetex
  \usepackage{unicode-math} % this also loads fontspec
  \defaultfontfeatures{Scale=MatchLowercase}
  \defaultfontfeatures[\rmfamily]{Ligatures=TeX,Scale=1}
\fi
\usepackage{lmodern}
\ifPDFTeX\else
  % xetex/luatex font selection
\fi
% Use upquote if available, for straight quotes in verbatim environments
\IfFileExists{upquote.sty}{\usepackage{upquote}}{}
\IfFileExists{microtype.sty}{% use microtype if available
  \usepackage[]{microtype}
  \UseMicrotypeSet[protrusion]{basicmath} % disable protrusion for tt fonts
}{}
\makeatletter
\@ifundefined{KOMAClassName}{% if non-KOMA class
  \IfFileExists{parskip.sty}{%
    \usepackage{parskip}
  }{% else
    \setlength{\parindent}{0pt}
    \setlength{\parskip}{6pt plus 2pt minus 1pt}}
}{% if KOMA class
  \KOMAoptions{parskip=half}}
\makeatother
\usepackage{xcolor}
\usepackage[margin=1in]{geometry}
\usepackage{color}
\usepackage{fancyvrb}
\newcommand{\VerbBar}{|}
\newcommand{\VERB}{\Verb[commandchars=\\\{\}]}
\DefineVerbatimEnvironment{Highlighting}{Verbatim}{commandchars=\\\{\}}
% Add ',fontsize=\small' for more characters per line
\usepackage{framed}
\definecolor{shadecolor}{RGB}{248,248,248}
\newenvironment{Shaded}{\begin{snugshade}}{\end{snugshade}}
\newcommand{\AlertTok}[1]{\textcolor[rgb]{0.94,0.16,0.16}{#1}}
\newcommand{\AnnotationTok}[1]{\textcolor[rgb]{0.56,0.35,0.01}{\textbf{\textit{#1}}}}
\newcommand{\AttributeTok}[1]{\textcolor[rgb]{0.13,0.29,0.53}{#1}}
\newcommand{\BaseNTok}[1]{\textcolor[rgb]{0.00,0.00,0.81}{#1}}
\newcommand{\BuiltInTok}[1]{#1}
\newcommand{\CharTok}[1]{\textcolor[rgb]{0.31,0.60,0.02}{#1}}
\newcommand{\CommentTok}[1]{\textcolor[rgb]{0.56,0.35,0.01}{\textit{#1}}}
\newcommand{\CommentVarTok}[1]{\textcolor[rgb]{0.56,0.35,0.01}{\textbf{\textit{#1}}}}
\newcommand{\ConstantTok}[1]{\textcolor[rgb]{0.56,0.35,0.01}{#1}}
\newcommand{\ControlFlowTok}[1]{\textcolor[rgb]{0.13,0.29,0.53}{\textbf{#1}}}
\newcommand{\DataTypeTok}[1]{\textcolor[rgb]{0.13,0.29,0.53}{#1}}
\newcommand{\DecValTok}[1]{\textcolor[rgb]{0.00,0.00,0.81}{#1}}
\newcommand{\DocumentationTok}[1]{\textcolor[rgb]{0.56,0.35,0.01}{\textbf{\textit{#1}}}}
\newcommand{\ErrorTok}[1]{\textcolor[rgb]{0.64,0.00,0.00}{\textbf{#1}}}
\newcommand{\ExtensionTok}[1]{#1}
\newcommand{\FloatTok}[1]{\textcolor[rgb]{0.00,0.00,0.81}{#1}}
\newcommand{\FunctionTok}[1]{\textcolor[rgb]{0.13,0.29,0.53}{\textbf{#1}}}
\newcommand{\ImportTok}[1]{#1}
\newcommand{\InformationTok}[1]{\textcolor[rgb]{0.56,0.35,0.01}{\textbf{\textit{#1}}}}
\newcommand{\KeywordTok}[1]{\textcolor[rgb]{0.13,0.29,0.53}{\textbf{#1}}}
\newcommand{\NormalTok}[1]{#1}
\newcommand{\OperatorTok}[1]{\textcolor[rgb]{0.81,0.36,0.00}{\textbf{#1}}}
\newcommand{\OtherTok}[1]{\textcolor[rgb]{0.56,0.35,0.01}{#1}}
\newcommand{\PreprocessorTok}[1]{\textcolor[rgb]{0.56,0.35,0.01}{\textit{#1}}}
\newcommand{\RegionMarkerTok}[1]{#1}
\newcommand{\SpecialCharTok}[1]{\textcolor[rgb]{0.81,0.36,0.00}{\textbf{#1}}}
\newcommand{\SpecialStringTok}[1]{\textcolor[rgb]{0.31,0.60,0.02}{#1}}
\newcommand{\StringTok}[1]{\textcolor[rgb]{0.31,0.60,0.02}{#1}}
\newcommand{\VariableTok}[1]{\textcolor[rgb]{0.00,0.00,0.00}{#1}}
\newcommand{\VerbatimStringTok}[1]{\textcolor[rgb]{0.31,0.60,0.02}{#1}}
\newcommand{\WarningTok}[1]{\textcolor[rgb]{0.56,0.35,0.01}{\textbf{\textit{#1}}}}
\usepackage{graphicx}
\makeatletter
\newsavebox\pandoc@box
\newcommand*\pandocbounded[1]{% scales image to fit in text height/width
  \sbox\pandoc@box{#1}%
  \Gscale@div\@tempa{\textheight}{\dimexpr\ht\pandoc@box+\dp\pandoc@box\relax}%
  \Gscale@div\@tempb{\linewidth}{\wd\pandoc@box}%
  \ifdim\@tempb\p@<\@tempa\p@\let\@tempa\@tempb\fi% select the smaller of both
  \ifdim\@tempa\p@<\p@\scalebox{\@tempa}{\usebox\pandoc@box}%
  \else\usebox{\pandoc@box}%
  \fi%
}
% Set default figure placement to htbp
\def\fps@figure{htbp}
\makeatother
\setlength{\emergencystretch}{3em} % prevent overfull lines
\providecommand{\tightlist}{%
  \setlength{\itemsep}{0pt}\setlength{\parskip}{0pt}}
\setcounter{secnumdepth}{-\maxdimen} % remove section numbering
\usepackage{bookmark}
\IfFileExists{xurl.sty}{\usepackage{xurl}}{} % add URL line breaks if available
\urlstyle{same}
\hypersetup{
  pdftitle={Homework 3},
  hidelinks,
  pdfcreator={LaTeX via pandoc}}

\title{Homework 3}
\author{}
\date{\vspace{-2.5em}}

\begin{document}
\maketitle

\textbf{This homework is due on Wednesday, October 22 2025 by 11:59pm.
Please show all your work and justify your answers. Use R to compute the
probabilities.}

\begin{enumerate}
\def\labelenumi{\arabic{enumi}.}
\tightlist
\item
  Trees are subjected to different levels of carbon dioxide atmosphere
  with 6\% of them in a minimal growth condition at 350 parts per
  million (ppm), 10\% at 450 ppm (slow growth), 47\% at 550 ppm
  (moderate growth), and 37\% at 650 ppm (rapid growth). What are the
  mean and standard deviation of the carbon dioxide atmosphere (in ppm)
  for these trees in ppm?
\end{enumerate}

\begin{Shaded}
\begin{Highlighting}[]
\CommentTok{\# minimal p=6\% 350ppm}
\CommentTok{\# slow p=10\% 450ppm}
\CommentTok{\# moderate p=47\% 550ppm}
\CommentTok{\# rapid p=37\% 650ppm}
\DocumentationTok{\#\# What are the mean and standard deviation of the carbon dioxide atmosphere for these trees in ppm? }
\NormalTok{rapid }\OtherTok{\textless{}{-}} \DecValTok{650}
\NormalTok{moderate }\OtherTok{\textless{}{-}} \DecValTok{550}
\NormalTok{slow }\OtherTok{\textless{}{-}} \DecValTok{450}
\NormalTok{minimal }\OtherTok{\textless{}{-}} \DecValTok{350}

\DocumentationTok{\#\# Mean calculation}
\CommentTok{\# 650 + 550 + 450 + 350 = 2000}
\CommentTok{\# 2000 / 4 = 500}

\CommentTok{\# create a vector}
\NormalTok{mean\_vector }\OtherTok{\textless{}{-}} \FunctionTok{c}\NormalTok{(rapid, moderate, slow, minimal)}

\CommentTok{\# calculate and print the mean}
\NormalTok{mean\_value }\OtherTok{\textless{}{-}} \FunctionTok{mean}\NormalTok{(mean\_vector)}
\CommentTok{\# print(sprintf("The mean is \%sppm", mean\_value))}



\DocumentationTok{\#\# Standard Deviation calculation}
\CommentTok{\# 1. find the mean  = 500}

\CommentTok{\# 2. subtract the mean from all the values}
  \CommentTok{\# 650 {-} 500 = 150}
  \CommentTok{\# 550 {-} 500 = 50 }
  \CommentTok{\# 450 {-} 500 = {-}50}
  \CommentTok{\# 350 {-} 500 = {-}150}
\NormalTok{a }\OtherTok{=}\NormalTok{ rapid }\SpecialCharTok{{-}}\NormalTok{ mean\_value}
\NormalTok{b }\OtherTok{=}\NormalTok{ moderate }\SpecialCharTok{{-}}\NormalTok{ mean\_value}
\NormalTok{c }\OtherTok{=}\NormalTok{ slow }\SpecialCharTok{{-}}\NormalTok{ mean\_value}
\NormalTok{d }\OtherTok{=}\NormalTok{ minimal }\SpecialCharTok{{-}}\NormalTok{ mean\_value}

\CommentTok{\# 3. Square the results}
  \CommentTok{\# 150 * 150 = 22500}
  \CommentTok{\# 50 * 50 = 2500}
  \CommentTok{\# ({-}50) * ({-}50) = 2500}
  \CommentTok{\# ({-}150) * ({-}150) = 22500}
\NormalTok{a\_2 }\OtherTok{=}\NormalTok{ a }\SpecialCharTok{*}\NormalTok{ a}
\NormalTok{b\_2 }\OtherTok{=}\NormalTok{ b }\SpecialCharTok{*}\NormalTok{ b}
\NormalTok{c\_2 }\OtherTok{=}\NormalTok{ c }\SpecialCharTok{*}\NormalTok{ c}
\NormalTok{d\_2 }\OtherTok{=}\NormalTok{ d }\SpecialCharTok{*}\NormalTok{ d}

\NormalTok{my\_variance\_vector }\OtherTok{\textless{}{-}} \FunctionTok{c}\NormalTok{(a\_2, b\_2, c\_2, d\_2)}
\CommentTok{\# print(my\_variance\_vector)}

\CommentTok{\# 4. Find the mean of the squared deviations (the variance)}
\NormalTok{my\_variance\_mean }\OtherTok{\textless{}{-}} \FunctionTok{mean}\NormalTok{(my\_variance\_vector)}
\CommentTok{\# print(my\_variance\_mean)}

\NormalTok{standard\_deviation }\OtherTok{\textless{}{-}} \FunctionTok{sqrt}\NormalTok{(my\_variance\_mean)}
\CommentTok{\# print(standard\_deviation)}



\FunctionTok{print}\NormalTok{(}\StringTok{"***** Final result *****"}\NormalTok{)}
\end{Highlighting}
\end{Shaded}

\begin{verbatim}
## [1] "***** Final result *****"
\end{verbatim}

\begin{Shaded}
\begin{Highlighting}[]
\FunctionTok{print}\NormalTok{(}\FunctionTok{sprintf}\NormalTok{(}\StringTok{"The Mean is \%sppm."}\NormalTok{, mean\_value))}
\end{Highlighting}
\end{Shaded}

\begin{verbatim}
## [1] "The Mean is 500ppm."
\end{verbatim}

\begin{Shaded}
\begin{Highlighting}[]
\FunctionTok{print}\NormalTok{(}\FunctionTok{sprintf}\NormalTok{(}\StringTok{"Standard Deviation is \%.1fppm."}\NormalTok{, }\SpecialCharTok{+}\NormalTok{ standard\_deviation))}
\end{Highlighting}
\end{Shaded}

\begin{verbatim}
## [1] "Standard Deviation is 111.8ppm."
\end{verbatim}

\begin{enumerate}
\def\labelenumi{\arabic{enumi}.}
\setcounter{enumi}{1}
\item
  Suppose that in the U.S. in 2011, 70\% of new start-ups failed in
  their first year. If there were 50 start-ups in the DMV area in 2011,
  what is the probability that 35 or more failed? What assumption did
  you need to make about the DMV?
\item
  A research team has developed a face recognition device to match
  photos in a database. From laboratory tests, the recognition accuracy
  is 92\% and trials are assumed to be independent.
\end{enumerate}

\begin{itemize}
\item
  If the research team continues to run laboratory tests, what is the
  mean number of trials until failure?
\item
  What is the probability that the first failure occurs on the tenth
  trial?
\item
  To improve the recognition algorithm, a chief engineer decides to
  collect 10 failures. How many trials are expected to be needed?
\end{itemize}

\begin{enumerate}
\def\labelenumi{\arabic{enumi}.}
\setcounter{enumi}{3}
\item
  You roll a die until obtaining two consecutive numbers that match.
  What it is the probability that it takes exactly \(x\) rolls to do
  this, for \(x\) an integer greater than 1?
\item
  The MIT-BIT arrhythmia database contains 48 heart signal recordings
  with 22 and 26 from female and male patients, respectively. For a
  classification task, an analyst randomly selects 36 records for
  predictive model training and keeps the other 12 records for testing
  the model performance.
\end{enumerate}

\begin{itemize}
\item
  What is the expected number of female patient recordings selected for
  training?
\item
  What is the probability that at least three male patient recordings
  are selected for training?
\item
  What is the probability that the same number of male and female
  patients are used for testing the model performance?
\end{itemize}

\begin{enumerate}
\def\labelenumi{\arabic{enumi}.}
\setcounter{enumi}{5}
\tightlist
\item
  An engineer is participating in a research project on the title
  patterns of junk emails. The number of junk emails which arrive in an
  individual's account every hour follows a Poisson distribution with a
  mean of 1.2.
\end{enumerate}

\begin{itemize}
\item
  What is the expected number of junk emails that an individual receives
  in an 8 hour day?
\item
  What is the probability that an individual receives more than two junk
  emails for the next three hours?
\item
  What is the probability that an individual receives no junk email for
  two hours?
\end{itemize}

\begin{enumerate}
\def\labelenumi{\arabic{enumi}.}
\setcounter{enumi}{6}
\tightlist
\item
  Consider a random variable \(X\) with probability density function
  \(f(x) = \theta x^{\theta-1}\) for \(0<x<1\), where \(\theta>0\) is
  the parameter that controls the pdf.
\end{enumerate}

\begin{itemize}
\item
  Show that \(f(x)\) is a valid pdf
\item
  Plot the pdf for \(\theta \in \{1/3,1/2,1,2,3\}\). Describe how the
  distribution changes for different values of \(\theta\).
\item
  Find a general expression for the expected value and variance of
  \(X\).
\item
  Find a general expression for the cumulative distribution function
  (cdf)
\end{itemize}

\begin{enumerate}
\def\labelenumi{\arabic{enumi}.}
\setcounter{enumi}{7}
\tightlist
\item
  Web crawlers need to estimate the frequency of changes to Web sites to
  maintain a current index for Web searches. Assume that the changes to
  a Web site follow a Poisson process with a mean of 6 days. Let a
  random variable \(X\) denote the time (in days) until the next change.
\end{enumerate}

\begin{itemize}
\item
  What is the probability that the next change occurs in less than 4.5
  days?
\item
  What is the probability that the time until the next change is greater
  9.5 days?
\item
  What is the time of the next change that is exceeded with probability
  90\%?
\end{itemize}

\end{document}
